\section{Laboratory and Precision Tests}

While astrophysical observations provide the tightest constraints on a 
possible photon mass, laboratory experiments play a crucial complementary 
role. They offer controlled, model-independent tests of electromagnetism 
over well-defined spatial scales and field configurations. Although 
laboratory limits are many orders of magnitude weaker than astrophysical 
ones, their conceptual significance is high: they test the foundations of 
Maxwell’s theory in environments where uncertainties about cosmic magnetic 
fields, plasma effects or large-scale structure do not enter.

\subsection*{Coulomb-law tests}

One of the most direct consequences of a non-zero photon mass is the 
modification of the Coulomb potential. Proca theory predicts a Yukawa-type 
potential,
\[
V(r) = \frac{q}{4\pi\varepsilon_0} \frac{e^{-m_\gamma r}}{r},
\]
instead of the familiar $1/r$ behaviour. Precision measurements of the 
electrostatic force therefore place bounds on the photon mass.

Laboratory tests using torsion balances, electrostatic null experiments and 
high-precision capacitor measurements typically yield
\[
m_\gamma \lesssim 10^{-14}\,\text{eV}.
\]
Although far weaker than astrophysical limits, these bounds rely only on 
basic laboratory physics and minimal theoretical assumptions.

\subsection*{Magnetic-field measurements}

A massive photon would imply that static magnetic fields decay over a 
finite range. Laboratory tests of long solenoids, superconducting magnetic 
shielding, and precision magnetometry have been used to search for such 
effects.

Experiments measuring the magnetic field inside superconducting cavities 
and around extended solenoid structures give constraints of the order
\[
m_\gamma \lesssim 10^{-13}\,\text{eV}.
\]
These tests are particularly clean because they rely on geometries in which 
magnetic field behaviour can be calculated with high precision.

\subsection*{Resonant-cavity experiments}

Resonant electromagnetic cavities provide another sensitive method for 
probing deviations from Maxwell’s equations. A photon mass would shift 
resonant frequencies and alter the mode structure of the cavity.

By comparing experimentally measured resonance spectra with theoretical 
predictions, one finds
\[
m_\gamma \lesssim 10^{-12}\,\text{eV}.
\]
Although these experiments are technologically demanding, they probe the 
wave nature of electromagnetism in a controlled and highly precise setting.

\subsection*{Interferometry and wave-propagation tests}

Interferometric techniques can detect tiny deviations in wave propagation. 
Even though laboratory baselines are short, modern interferometers achieve 
sensitivities that allow for stringent null tests of Maxwellian behaviour.

Experiments using:

\begin{itemize}
	\item microwave interferometry,
	\item optical delay lines,
	\item and laser cavities
\end{itemize}

have all searched for photon-mass–induced dispersion or phase anomalies.  
These experiments consistently support Maxwell’s theory and set bounds at 
the level of
\[
m_\gamma \lesssim 10^{-13}\,\text{eV}.
\]

\subsection*{Comparison with astrophysical bounds}

Laboratory tests cannot compete with astrophysical constraints in absolute 
sensitivity. However, they provide a crucial independent foundation:

\begin{itemize}
	\item They rely on simple, well-controlled geometries.
	\item They are free from assumptions about cosmic magnetic fields or 
	plasma densities.
	\item They test the structure of electromagnetism under reproducible 
	conditions.
	\item Their null results are fully consistent with the massless-photon 
	predictions of Maxwell and QED.
\end{itemize}

The agreement between laboratory and astrophysical bounds — despite their 
very different physical regimes — strongly reinforces the conclusion that 
the photon is, to extremely high precision, massless.

\subsection*{Transition to conceptual arguments from QED}

The next chapter turns from empirical tests to theoretical consistency 
requirements. Within quantum electrodynamics, the masslessness of the 
photon is intimately connected to gauge symmetry, renormalizability, and 
the structure of quantum amplitudes. These considerations provide the 
deepest conceptual reasons why a photon mass is incompatible with the 
Standard Model.

Laboratory tests therefore complete the empirical foundation of the 
argument: if the photon has a mass at all, it must be far too small to be 
detected by any current experiment.
